\section{模型架构与强化学习的适配}

\subsection{算法适配架构 vs. 架构适配算法}

崔干曲提出了一个有趣的视角转换：通常我们考虑的是"让算法适应模型架构"，但很少有人反过来想"让架构适应算法"。他认为架构设计更多受到Infra和硬件效率的驱动（训练效率、推理效率、多模态支持、长文本处理），RL算法处于相对"靠后"的位置——在架构确定之后再进行适配。

\begin{warningbox}{Hybrid架构带来的新挑战}
在Hybrid架构（如Mamba + Attention混合模型）上做RL训练，复杂度远高于Full Attention架构。Hybrid架构涉及State的更改和回滚操作，高效计算这些操作并非trivial。这也意味着RL训练与架构设计之间存在Trade-off——RL研究者需要与架构团队协调，确保新架构对后续RL训练友好。
\end{warningbox}

\subsection{RL友好即推理友好}

胡建从infra角度给出了一个精辟的判断：

\begin{importantbox}{对RL友好的模型架构 = 对推理友好的模型架构}
在当前Long Context、Long Thought输出的场景下，RL训练的大部分算力瓶颈集中在生成（Generation/Rollout）部分。因此，提升RL训练效率本质上就是提升推理效率。推理友好的架构创新包括：
\begin{itemize}
    \item \textbf{Linear Attention / Mamba}：降低长序列的计算复杂度
    \item \textbf{千问3 Next}：在推理友好性方面的架构创新
    \item \textbf{MTP（Multi-Token Prediction）}：DeepSeek提出，对推理效率有显著提升
\end{itemize}
\end{importantbox}

胡建同时提出了一个前瞻性观点：如果未来\textbf{Diffusion LLM}取代了Autoregressive范式，整个RL训练的瓶颈可能从推理优化转向训练优化——这将是一个"天翻地覆"的变化。

\subsection{两种看待RL与大模型关系的视角}

郑楚杰提出了两种对立的视角：

\textbf{视角一：RL是大模型的辅助。}先确定模型架构（MoE、MTP等推理加速技巧），再在此基础上适配RL算法。这是当前工业界的主流做法。

\textbf{视角二：大模型是RL的辅助。}如果我们的终极目标是通过RL实现AGI，那么大模型只是提供先验知识的手段。在这种视角下，应该使用最有利于RL训练的模型架构（如Dense Model，不加MTP等加速技巧）。但这种做法在实践中因成本太高而少有人采用。

\begin{knowledgebox}{模型架构变化影响RL特性}
郑楚杰的一个重要观察：模型架构（以及基模预训练）的变化会连带改变后续RL训练的特性。因此，RL问题和架构问题不能分开来看——每一代模型都需要针对其特有的问题做RL方面的改进。
\end{knowledgebox}

\subsection{MoE架构的RL训练困境}

李英儒从优化理论的角度解释了为什么MoE模型难以进行RL训练：

\begin{warningbox}{MoE难训的理论根源}
\begin{itemize}
    \item MoE架构的离散选择性（Router的Expert选择）从根本上就给优化带来了困难
    \item 从Off-Policy算法的角度看，MoE实际上\textbf{缩小了Trust Region}——Expert的选择变化会导致模型行为的剧烈变化
    \item 现有的Routing Replay等技巧都是在解决这一Trust Region缩小的问题
    \item 这个问题不是RL独有的——MoE在Pretraining中就已经遇到类似的路由平衡问题
\end{itemize}
\end{warningbox}

李英儒对未来持乐观态度：可能会出现一些新的Sparse架构，既能保持Sparse激活的效率优势，又能在优化上具有更好的特性（例如，在Router上更好地近似梯度）。

\subsection{本章小结}

模型架构与RL训练之间存在深层的耦合关系。当前工业界以"架构优先，RL适配"为主流范式。RL友好性本质上等价于推理友好性——架构层面的创新（Linear Attention、MTP等）直接惠及RL训练效率。MoE架构给RL带来了独特的Trust Region缩小问题，需要算法与架构协同设计来解决。


